% TOP_PROBLEM: {"id": 3100079, "title": "There is (essentially) a unique sequence over {1,2} which is its own run-length encoding", "category_id": 2, "category": {"id": 2, "name": "combinatorics", "display_name": "Combinatorics", "description": "Counting problems, graph theory, discrete structures.", "slug": "combinatorics", "order_index": 2, "created_at": "2026-07-31T15:26:25.671Z"}, "status": "partially_solved", "problem_number": "AMR-030-0079", "set_id": 13}
% FIRST_POSTED: 2026-10-02
% ENTRY_KIND: statement_only
% UnsolvedMath ID 3100079; source catalogue code AMR-030-0079
% !TeX program = lualatex
% Definition and sources only; this file is not an LLM solution attempt.
\documentclass[11pt,a4paper]{article}
\usepackage[margin=25mm]{geometry}
\usepackage{fontspec}
\setmainfont{lmroman10-regular.otf}[BoldFont=lmroman10-bold.otf,
 ItalicFont=lmroman10-italic.otf,BoldItalicFont=lmroman10-bolditalic.otf]
\usepackage{amsmath,amssymb,mathtools}
\usepackage{unicode-math}
\setmathfont{latinmodern-math.otf}
\AtBeginDocument{\let\setminus\smallsetminus}
\usepackage{xurl,hyperref}
\hypersetup{colorlinks=true,urlcolor=blue}
\setcounter{secnumdepth}{0}
\setlength{\parindent}{0pt}
\setlength{\parskip}{5pt}
\setlength{\emergencystretch}{3em}
\begin{document}
\section*{There is (essentially) a unique sequence over \{1,2\} which is its own run-length encoding}\label{3100079}
{\small UnsolvedMath ID 3100079 · AMR-030-0079 · Combinatorics · AMR Open Problem Lists}
\subsection{Definitions and mathematical statement}
A run in an infinite word is a maximal consecutive block of
equal symbols. If all runs are finite, its run-length encoding
is the sequence of their lengths, listed from left to right.
The standard Kolakoski sequence is the infinite word
$K=(K_1,K_2,\ldots)$ over $\{1,2\}$ beginning with $1$ whose
run-length encoding equals $K$ itself. Its initial symbols are
\[
                 1,2,2,1,1,2,1,2,2,1,2,2,\ldots.
\]
Equivalently, the successive runs alternate between symbols
$1$ and $2$ and have successive lengths $K_1,K_2,\ldots$;
the initial-symbol convention specifies the standard sequence.

Does its natural density of ones exist and equal one half?
In explicit terms, is
\[
           \lim_{N\to\infty}\frac1N
                  \bigl|\{1\le j\le N:K_j=1\}\bigr|=\frac12\,?
\]
Both existence of the limit and its value are part of the
question. A proof of equality of some upper or lower densities,
or of frequencies along selected subsequences of prefix lengths,
would require additional work to establish the displayed limit.

The count is by positions in the original word, not by runs.
Runs alternate symbols, so their frequencies by run number are
automatically balanced; their variable lengths are precisely
what makes the frequency of symbols a separate issue.
\subsection{Short English statement}
Does the standard Kolakoski sequence have a natural density of ones equal to one half?
\subsection{Sources}
\begin{itemize}
\item[S1] ULAM.AI and the UnsolvedMath Contributors, supplied problems.json, record 3100079 (AMR-030-0079), AMR Open Problem Lists. Catalogue identity and source transcription; CC BY 4.0.
\url{https://huggingface.co/datasets/ulamai/UnsolvedMath}.
\item[S2] Cooper - Combinatorial Problems I Like (2020 snapshot), Words and Codes, source-order bullet 79. Original source indexed by AMR-030-0079.
\url{https://people.math.sc.edu/cooper/combprob.html}.
\end{itemize}
\end{document}
